\subsection{The comparison after the saved blind proof}

\begin{statement}{P5.7}{Known}
The independently derived values in P5.3--P5.6 agree with the corresponding
statements in version~1 of Du--Guo--Kusakabe--Xie,
\emph{Oka Geometry and Deformations of Torus-Fibred Complex Six-Spheres}.
Version~2, titled \emph{Oka Geometry of Complex Six-Spheres}, does not
contain those tangent-cohomology, miniversality, full-automorphism or
exact-period assertions.
\end{statement}
\begin{proof}
Checkpoint~01 was created at 2026-11-17 00:49:44 UTC and its persistent
save succeeded before 00:50:09 UTC. It contains both complete blind
proof files. Their SHA-256 hashes and the successful save are recorded
in \path{records/checkpoint_01.json} and
\path{records/checkpoint_01_saved.json}. The first attempts to open the
comparison paper occurred subsequently, at 00:50:09 UTC. Those web
requests failed; direct version-specific PDF downloads succeeded at
00:50:32 UTC. We read the downloaded statements and proofs only after
that point. Thus the comparison did not enter the derivation above.

In version~1, Theorem~6.12, printed pp.~81--82, gives the same four
tangent dimensions as P5.3. Propositions~6.15 and~6.16, pp.~84--85,
give the nonzero Kodaira--Spencer map and the smooth one-dimensional
miniversal germ of P5.4. Theorem~7.14, p.~95, gives the full group of
P5.5. Theorem~1.4, p.~3, gives the period relation of P5.6. These
identifications use the same additive period entry; replacing the
particular solution by another adds a constant to both parameters
and leaves their difference unchanged. This supplies the parameter
comparison, rather than an identification inferred from notation.

The direct-image bundles in version~1 Proposition~6.10 coincide with
those in P5.2 after identifying its orbifold line
$\mathcal O_{\widehat B}(-1)$ with our $\widehat L$: both have degree
$-1/12$ and characters $2$ modulo~3 and $1$ modulo~4. Their quotient
descends to a degree-zero line on $\mathbb P^1$, hence is trivial.
The proof above computes ordinary tangent cohomology through the
critical ideal $\mathcal I_D$; version~1 passes through the
logarithmic tangent bundle and the normal line of the normalization.
Both quotient sequences are obtained from the same local equations
$q=xy$ and $q=xyz$, and both give the displayed dimensions.

We read version~2's introduction and complete contents and searched
its full text for the four topics. Its body has Sections~1--4 and
Appendices~A--B; the former Sections~6--7 are absent. Consequently
version~2 supplies no competing value for these four controls.
The removal alone supplies no mathematical refutation. Our proofs
remain P5.1--P5.6. The recorded literature status is an independent
proof of results already stated and proved in version~1; no priority
is claimed. Both exact downloaded versions and the access log are
retained under \path{records/part5_comparison/}.
\end{proof}
